% ============================================================
\section{Limitations, Open Problems, and Regulatory Framing}
\label{sec:limitations}

M\textsuperscript{2}'s structural guarantees are conditional. Section~\ref{sec:protocol-threat-model} delimited the threat model; this section enumerates the residual risks that lie outside it, identifies open mathematical problems, and offers an academic regulatory framing for readers unfamiliar with the digital-asset compliance landscape. We follow the FC convention of separating ``Limitations'' (axes along which the contributions are bounded) from ``Open Problems'' (questions on which we have made partial progress but cannot fully answer in this paper).

\subsection{Smart-contract residual risk}
\label{sec:lim-contract}

The protocol is immutable. There is no upgrade path, no admin pause, and no governance mechanism that could deploy a bug fix. The mitigation is a standard audit-and-formal-verification stack: external code review, an internal differential-testing suite against a reference implementation, and the Halmos / Certora mechanized-verification path of Section~\ref{sec:formal-mechanized}. None of these can eliminate residual risk. A contract whose proof of correctness rests on a paper proof and an auditor's reading is more trustworthy than one without those checks but is not bug-free. We make no claim of zero residual contract risk, only of a small one whose magnitude we have made every effort to bound.

\subsection{Backing-stable issuer risk}
\label{sec:lim-stable}

The treasury holds a single deployer-chosen USD-pegged ERC-20. The protocol's invariants hold regardless of the backing stable's market price (Theorem~\ref{thm:floor-monotonicity} is in units of the backing stable, not dollars). They do not hold if the issuer freezes the treasury balance: the March 2023 USDC depeg, in which Circle's exposure to Silicon Valley Bank produced a $\sim 13\%$ intra-day deviation from the dollar peg, did not freeze any specific address, but issuer-side freeze authority on USDC, USDT, and most centralized stables remains a residual capability that the protocol cannot defend against. Mitigations available to deployers: (i) choose a stable whose freeze authority is structurally constrained (e.g., a fully algorithmic stable with no centralized issuer---at the cost of accepting that stable's depeg risk); (ii) accept this risk and disclose. The protocol is structurally single-stable; diversification would require an entirely different design.

\paragraph{Open problem (stable-peg sensitivity).} Characterize the effective floor as a function of the backing stable's mark-to-dollar price under various depeg dynamics. This is straightforward for parametric depeg models but non-trivial under regime-switching or correlation with broader risk events; we leave it to follow-up work.

\subsection{Operator-revenue dependence}
\label{sec:lim-operator}

The floor's growth (not its existence) depends on continued revenue deposits by the operator. Theorem~\ref{thm:spot-floor-convergence} shows the system gracefully degrades to a CEF analog at zero revenue---the floor stops compounding but does not retreat. Operator fraud (revenue collected but routed elsewhere) is not in the protocol's threat model: this is a real-world failure mode addressed by entity-level disclosure, audit committee oversight, and qualified legal counsel, not by the on-chain protocol.

\paragraph{Open problem (operator IC constraint).} Model the operator as a strategic agent with an outside option (e.g., dumping vested tokens versus depositing revenue) and derive the incentive-compatibility constraint as a function of the operator's vested-token holdings, the vesting cliff schedule, and the time horizon. The IC constraint conditions when continued deposits are a dominant strategy for the operator, formalizing the qualitative intuition that vesting + long cliff align operator and holder interests.

\subsection{LP exhaustion as a feature}
\label{sec:lim-lp}

The LP becomes thin under sustained buy-and-burn (Theorem~\ref{thm:lp-half-life}). Spot prices diverge from floor; slippage on large swaps becomes punitive. The protocol does not refill the LP. We argue this is an intentional feature, not a bug: any refill mechanism would either (i) require post-genesis minting (violating the supply-only-decreases invariant) or (ii) consume treasury stable to redeposit into the LP, which strictly lowers the floor as $\Floor_{\text{new}} = (\Tres - X)/(\Sup + Y) < \Tres/\Sup$ for any positive injection $X$ paired with new tokens $Y$. Both alternatives are inconsistent with the load-bearing properties of the design. The terminal ``thin LP, redemption-dominant'' regime is the system's intended steady state; it is not a failure.

\subsection{Open mathematical problems}
\label{sec:lim-open-math}

We identify seven open problems whose resolution would strengthen the paper but lies beyond its scope.

\paragraph{(P1) Dynamic fees as a function of LP depth.} The 0.10\%/3\% fee asymmetry is constant. A state-dependent fee schedule (e.g., $f_s$ scaling with the spot-floor gap) could improve the run-resistance properties of Theorem~\ref{thm:bank-run}. Designing such a schedule that preserves the closed-form $A^*$ analysis is an open question.

\paragraph{(P2) Multi-block MEV under reorg.} Our MEV analysis (Theorem~\ref{thm:mev-bound}) assumes single-block adversary windows. Multi-block MEV under builder collusion on PoS chains (\citet{Daian2020}'s ``Flash Boys 3.0'' regime) is a known unresolved problem class for any AMM design; we explicitly do not solve it here.

\paragraph{(P3) Oracle-free spot manipulation.} M\textsuperscript{2} does not consume an external oracle, but external protocols may consume M\textsuperscript{2}'s spot as a price oracle. Manipulating spot to fool such consumers is a third-party-externality problem we do not address.

\paragraph{(P4) Cross-block ordering and MEV-relay collusion.} Builder-relay collusion can produce cross-block ordering effects that violate the single-block independence assumption in our MEV bound. The bound degrades gracefully but is not tight under this adversary.

\paragraph{(P5) End-to-end mechanized verification.} We have given a paper proof of Theorem~\ref{thm:floor-monotonicity}. Discharging it on the deployed bytecode under Halmos or Certora is in progress and will be reported in follow-up work.

\paragraph{(P6) Quantitative LVR/floor-capture identity.} Conjecture~\ref{conj:lvr-capture} states a directional inversion of the standard LVR framing of \citet{MilionisMoallemiRoughgarden2023}: for the protocol-owned LP position, LVR-equivalent value is captured into the floor rather than lost to arbitrageurs. Promoting this to a quantitative theorem requires a unit reconciliation between LVR's per-unit-time stochastic-calculus formulation and M\textsuperscript{2}'s discrete event-driven floor framing, together with a careful per-leg accounting of the asymmetric fee. We have not completed this reconciliation; the directional content is sufficient for the qualitative arguments in Sections~\ref{sec:related-formal} and~\ref{sec:formal-lvr}, but the quantitative identity is left open.

\paragraph{(P7) Sequenced bank-run strategies.} Theorem~\ref{thm:bank-run} bounds the attacker's proceeds within the \emph{simultaneous}-strategy class, in which the dump and the redemption occur in a single block without an intervening \texttt{collectFees} realization. A sequenced strategy that interleaves LP dumps, \texttt{collectFees} calls, and redemptions within or across blocks can extract additional value scaling with the attacker's residual holdings (Section~\ref{sec:econ-bank-run}, ``Sequenced strategies''). Characterizing the supremum of attacker proceeds across the full sequenced-strategy class---which requires modeling the equilibrium response of the \texttt{collectFees}-bounty payer and the v4 fee-realization timing dynamics, as well as the cross-block adversary's interaction with the (P4) builder-relay-collusion class---is left open.

\subsection{Regulatory framing for academic readers}
\label{sec:lim-regulatory}

We close with a brief academic positioning on the regulatory question. This is not legal advice; any deployer must obtain qualified counsel in every relevant jurisdiction (operator domicile, target user geographies, listing venues). Our purpose here is to map M\textsuperscript{2}'s on-chain properties to the dimensions along which existing securities and digital-asset frameworks evaluate tokens.

\paragraph{Securities-style frameworks.} Howey-style tests (U.S.\ securities law), the Markets in Crypto-Assets (MiCA) regulation (EU), the Virtual Asset Regulatory Authority (VARA) rules (Dubai), and the Monetary Authority of Singapore (MAS) Payment Services Act each ask different versions of the question ``does this asset represent an investment contract.'' M\textsuperscript{2} has features that weigh in both directions: the redemption-guarantee and the operator's affirmative duty to route revenue could be read as a fixed-return commitment (weighing toward securities classification); the absence of any pooled-investment management vehicle, the floating spot, the immutability of the parameters, and the protocol's structural independence from the operator's post-genesis discretion weigh against. We take no position on the outcome in any jurisdiction. Counsel will produce a different answer in each.

\paragraph{Spectrum positioning.} On the canonical asset spectrum (utility token---governance token---asset-backed---e-money---security), M\textsuperscript{2} is closest to ``asset-backed'' in the sense of an asset with a contractual claim on a reserve. It is explicitly not an e-money token under MiCA (no peg, no redemption-at-par); it is not a pure utility token (the floor is the central feature); it has no governance, so it is not a governance token. The closest TradFi categorization is ``asset-backed warrant or note'' or, in some jurisdictions, a redeemable preferred-share analog---each of which is treated differently in different regimes.

\paragraph{Jurisdictional reach.} Smart contracts are address-agnostic and globally reachable. Geofencing operates at the front-end (user-interface) layer, not at the contract layer. A deployer who wishes to restrict access in specific jurisdictions can do so only through interface-level filters, and these can be bypassed by direct contract interaction. This is a feature of the substrate, not of M\textsuperscript{2} specifically, and the legal exposure it produces should be evaluated by counsel for every target jurisdiction.

\paragraph{Recommendation.} Deployers should obtain counsel before launch. Where the public materials describing the protocol use the language of ``guaranteed floor'' or ``redemption,'' counsel should review them; while these are accurate technical descriptions of contract behavior, regulators across jurisdictions read them as making investment-contract claims. This paper makes the technical claims; the legal characterization is a separate exercise.
